\section{Detailed Analysis}

\subsection{Layer-Wise Sensitivity and Error Dynamics}
Reconstruction error profiles across transformer layers reveal significant heterogeneity. Early projection layers exhibit higher sensitivity to rank truncation and aggressive pruning due to initial low-level feature encoding requirements. In contrast, intermediate layers display steep singular value decay, tolerating aggressive compression budgets without functional degradation.

\begin{table}[H]
\centering
\small
\begin{tabular}{lcccc}
\toprule
\textbf{Layer Group} & \textbf{Projections} & \textbf{Spectral Decay Rate} & \textbf{Optimal Rank} & \textbf{Reconstruction Error} \\
\midrule
Early (Layers 0--2) & Q, K, V, O & Slow ($\alpha = 0.85$) & $r=192$ & $0.182$ \\
Middle (Layers 3--8) & Q, K, V, O & Rapid ($\alpha = 1.42$) & $r=96$ & $0.114$ \\
Late (Layers 9--11) & Q, K, V, O & Moderate ($\alpha = 1.08$) & $r=128$ & $0.145$ \\
\bottomrule
\end{tabular}
\caption{Layer-wise spectral sensitivity across transformer depth.}
\label{tab:layer_sensitivity}
\end{table}

\subsection{Memory Ceilings and Accelerator Scaling}
Memory profiling indicates that \texttt{LowRankLinear} prevents GPU memory spikes by eliminating redundant intermediate projection buffers. Bounded memory consumption facilitates higher batch concurrency during server-side deployment. Under sequence lengths $L \in \{128, 256, 512\}$, peak VRAM footprint remains within the bounds of standard edge hardware.

\subsection{Ablation: Impact of Residual Rank on Recovery Trajectories}
We evaluate the influence of adapter rank $r_L \in \{4, 8, 16, 32, 64\}$ on downstream SST-2 validation accuracy recovery:
\begin{table}[H]
\centering
\small
\begin{tabular}{lcccc}
\toprule
\textbf{Adapter Rank $r_L$} & \textbf{Trainable Params} & \textbf{Accuracy (\%)} & \textbf{Recovery Time (s)} & \textbf{VRAM (MB)} \\
\midrule
$r_L = 0$ (Frozen) & 0 & 81.2 & 0 & 124.0 \\
$r_L = 4$ & 147,456 & 86.4 & 18.2 & 128.5 \\
$r_L = 8$ & 294,912 & 87.8 & 24.5 & 132.0 \\
$r_L = 16$ (Default) & 589,824 & 88.5 & 35.8 & 142.4 \\
$r_L = 32$ & 1,179,648 & 88.7 & 54.2 & 168.0 \\
\bottomrule
\end{tabular}
\caption{Ablation of recovery adapter rank $r_L$ on DistilBERT.}
\label{tab:ablation_adapter_rank}
\end{table}

\subsection{Subspace Preservation under Adaptation}
The empirical stability observed post-adaptation confirms that the low-rank basis identified by \texttt{soka} effectively spans the principal representation subspace.
